\section{Rebuttal}
We thank the reviewers for their detailed feedback and valuable suggestions to help us further improve this work. We address their concerns through this response and will include more detailed results in the final version.

\subsection{Response to Review-1:}

\paragraph{(1) Results of Calibration:}
Our method has demonstrated strong performance on the selective classification, which is directly related to model calibration and requires an accurate estimation of the total uncertainty. As per the reviewer's suggestion, we computed the ECE scores of our methods and other baseline methods. 
CIFAR10: Ours(10.4), Base-Model(13.2), LULA(10.8), BeBayesian(10.6), MCDropout(11.3); 
CIFAR100: Ours(7.3), Base-Model(26.2), LULA(7.1), BeBayesian(5.8), MCDropout(22.5). 
It shows that our method achieves comparable or better calibration performance than the baseline methods.

\paragraph{(2) Comparison with existing meta-modeling approach:}
To the best of the authors' understanding, Chen et al. 2019 and Jain et al. 2021 are the only existing meta-model-based UQ works. It’s possible that we misunderstood the reviewer’s request, but we already compared with Chen et al. 2019 (denoted as "Whitebox") in the experiment section and demonstrated that our method outperformed it across different tasks. Additionally, we note that Jain et al. 2021 mainly focuses on regression instead of classification and should be considered contemporaneous. We leave comparison with this work for future work.
%, its problem setting mainly focuses on regression instead classification
%, and it is still unpublished. Thus, a direct comparison may not serve the purpose at this time.

\subsection{Response to Review-2:}

\paragraph{(1)}
Thanks to the reviewer for pointing out the typo in Table 1. We will correct it in the revision.

% , and we are sure that our experiment setups are all accurate

\paragraph{(2) Difference between our work and prior network:}

First, the motivations and problem settings are different. Our work tries to improve the uncertainty quantification performance of an arbitrary pre-trained model in a post-hoc manner. Prior Network [1] aims to train a Bayesian model with the UQ capability from scratch. Second, the Dirichlet technique used in this work, though it might be similar to [1], simply serves as one efficient technique to quantify epistemic uncertainty. However, our method is not limited to using the Dirichlet technique as long as the output of the meta-model can quantify total and epistemic uncertainties. Moreover, compared to [1], our proposed method uses the variational loss instead of the KL divergence loss that requires additional OOD data defined by Equation 12 in [1]. In the revision, we will further clarify the main differences.

\subsection{Response to Review-5:}

\paragraph{(1) Comparison with existing OOD work ALOE[1]:}

The relation between UQ and OOD has been discussed in section 5.1, and we agree that we can introduce more about OOD in the related work. The scope of this work is to propose a general approach that can quantify different types of uncertainties for various tasks, especially in a challenging post-hoc setting. The OOD detection task considered here is a specific application that highly relies on good epistemic uncertainty. However, the existing OOD detection approach is not necessarily capable of the general UQ capability for other applications, such as selective classification that relies on total uncertainty. 


As the reviewer suggested, we compare with [1] to further validate our method.  The results are the AUROC scores evaluated on CIFAR 10 and CIFAR 100 (average over five OOD datasets). CIFAR10: Ours-MI(98.4), ALOE(97.3);  CIFAR100: Ours-MI(87.4), ALOE(86.9). Our method can still demonstrate advantages over ALOE on OOD tasks. Moreover, ALOE [1] requires an additional OOD dataset for training their model, while our method does not. Besides, we notice that training ALOE is computationally expensive. %that takes half day.
We will cite [1] in the final version and provide more discussion and comparison. 

\paragraph{(2) Comparison with Calibration approach [2]:}

Traditional calibration approaches, such as temperature scaling [2], have two major limitations. First, it requires an additional validation dataset to optimize the temperature. Second, as we mentioned in the third paragraph of the introduction section, calibration methods fail to capture epistemic uncertainty, which is crucial for OOD detection. To further elaborate this argument, we evaluate the performance of the calibration approach [2] on OOD detection and provide the average AUROC score: CIFAR10: Ours-MI(98.4), Base-Model(87.9), Calibration[2](92.3); CIFAR100: Ours-MI(87.4), Base-Model(75.8), Calibration[2](82.5). The calibration can improve the performance over the base model, but our proposed method can still outperform calibration [2] by a large margin.


\subsection{Response to Review-6:}

\paragraph{(1) Demonstrate the efficiency:}

Our meta-model-based approach is much more efficient than traditional uncertainty quantification approaches because it has a simpler structure and faster convergence speed. To quantitatively demonstrate such efficiency, we measure the wall-clock time of training the meta-model in seconds (on a single Titan RTX GPU): CIFAR10/VGG16: 66.5s (for five epochs); CIFAR100/WideResNet: 241.9s (for ten epochs). The training time of the meta-model can be negligible compared to those approaches training the entire base model from scratch (usually taking several hours).








%As the reviewer suggested, we compare with [1] to further validate our method.  The results are the AUROC scores evaluated on CIFAR 10 and CIFAR 100 (average over five OOD datasets). CIFAR10: Ours-MI 98.4(100.0, 98.8, 95.2, 98.1, 100.0), ALOE 97.3(99.9, 98.3, 92.1, 99.9, 96.2); CIFAR100: Ours-MI 87.4(94.3, 84.4, 81.9, 85.5, 90.8), ALOE 86.9(99.9, 86.9, 70.2, 85.3, 92.2). Our method can still demonstrate advantages over ALOE on OOD tasks. Moreover, ALOE [1] requires an additional OOD dataset for training their model, while our method does not. Besides, we notice that training ALOE is computationally heavy that takes half day. We will cite [1] in the final version and provide more discussion and comparison. 